% ------------------------------------------------------------------------
% bjourdoc.tex for birkjour.cls*******************************************
% ------------------------------------------------------------------------
%%%%%%%%%%%%%%%%%%%%%%%%%%%%%%%%%%%%%%%%%%%%%%%%%%%%%%%%%%%%%%%%%%%%%%%%%%

\documentclass{birkjour}
%
%
\usepackage{amssymb}
\usepackage{hyperref}
% THEOREM Environments (Examples)-----------------------------------------
%
 \newtheorem{thm}{Theorem}[section]
 \newtheorem{cor}[thm]{Corollary}
 \newtheorem{lem}[thm]{Lemma}
 \newtheorem{prop}[thm]{Proposition}
 \theoremstyle{definition}
 \newtheorem{defn}[thm]{Definition}
 \theoremstyle{remark}
 \newtheorem{rem}[thm]{Remark}
 \newtheorem*{ex}{Example}
 \numberwithin{equation}{section}



%%%%%%%%%%%%%%%%%%%%%%%
\usepackage{hyperref}
%%%%%%%%%%%%%%%%%%%%%%%%
\author[S. T. Ura]{Sergio Tsuyoshi Ura}
\email{sergioura@gmail.com}


\author[N. C. L. Penteado]{Northon Canevari Leme Penteado}
\address{Instituto de Formação de Educadores
\\
Universidade Federal do Cariri
\\
CEP 63260-000- Brejo Santo- CE
\\
Brazil}
\email{northon.canevari@ufca.edu.br}

\author[Vituri]{Guilherme Vituri}

\author[Mio]{Washington Mio}

\author[alice]{Alice Kimie Miwa Libardi}



\keywords{Topological Data Analysis, Persistent homology, Diptera classification, Wing venation}


\begin{document}



\title[Diptera wing classification using Topological Data Analysis]{Diptera wing classification using Topological Data Analysis}


\begin{abstract}
We apply tools from Topological Data Analysis (TDA) to classify Diptera families based on wing venation patterns. Using three complementary filtration strategies: Vietoris-Rips on point clouds, radial filtrations, and directional (height) filtrations on wing images. We extract \(H_0\) and \(H_1\) topological features via extended summary statistics and compare classifiers via leave-one-out cross-validation. We focus on interpretable models (LDA, Decision Trees) to identify explainable topological criteria that distinguish families.


\end{abstract}

%%% ----------------------------------------------------------------------
\maketitle
%%% ----------------------------------------------------------------------
{\bf DIRECIONAMENTO PARA O VITURI: 
\\
1) escrevemos este texto para que sugerimos par o artigo. Alteramos aa ordem do que está na última versão \url{https://juliatda.github.io/tda-fly/#feature-ablation} 
\\
2) está em destaque no texto onde colocar as imagens, de acordo com texto citado acima. No decorrer do texto tem por exemplo ***figure 3.1, ai no texto \url{https://juliatda.github.io/tda-fly/#feature-ablation}  terá a imaagem na seção 3.1.
\\
3) as imagens tem que ajustar os tamanhos, pois ficaram grandes demais.
\\
4) a ordem dos autores, como está neste teex, foi sugestão da Alice
\\
5) a introdução e as referências estão no último arquivo que alterei, na branch que criamos para eu mexer no github 
}

\section{Introduction}
...

\section{Methods}

 introdução

\subsection{Data loading and preprocessing}\quad

In the analysis we use 70 images of dipteras wings with 9 families: "Asilidae", "Bibionidae", "Ceratopogonidae", "Chironomidae", "Rhagionidae", "Sciaridae", "Simuliidae", "Tabanidae", "Tipulidae".

{\noindent Samples per family:

  Asilidae: 8
  
  Bibionidae: 6
  
  Ceratopogonidae: 8
  
  Chironomidae: 8
  
  Rhagionidae: 4
  
  Sciaridae: 6
  
  Simuliidae: 7
  
  Tabanidae: 11
  
  Tipulidae: 12}

{\bf COLOCAR IMAGENS DE TODAS AS ASAS}


These images are loading, binarized and processing applying the following techniques:

- Gaussian blur: to close small gaps in the wings membranes and keep it connected.

- Crop to the bounding box: to standardize the sizes of the images. 

- Connect pixel components: some images contain noise points that is not good to analysis to eliminate the noise. This process became the images into a single connected component, prioritizing those with the largest number of disconnected components before correction.

{\bf COLOCAR ALGUMAS IMAGENS COMPARANDO BASTANTE COMPONENTES CONEXAS E DEPOIS DO PROCESSO SOMENTE UMA COMPONENTE E A TABELA COM OS DADOS DO PROCESSO}

We selected 750 pixels of each image to do a Topological Data Analysis.



\subsection{Filtrations}
\quad


The \textbf{persistent homology} is the main tool of Topological Data Analysis (TDA), it is a  method for computing topological features at different levels.
To compute the persistent homology of space we must obtain a filtration of  a simplicial complex that represent the space. 



In this project, we compute persistent homology using three filtration strategies. 


\subsubsection{Strategy 1: Vietoris-Rips filtration on point clouds}
\quad

Given a set of points in \(\mathbb{R}^{n}\), the Vietoris-Rips complex at scale \(\epsilon\) connects any subset of points that are pairwise within distance \(\epsilon\). 
As \(\epsilon\) increases from \(0\), we obtain a nested sequence of simplicial complexes: the Rips filtration. 
This is the most common filtration in TDA for point-cloud data, but it is computationally expensive (since it must consider all pairwise distances), which is why we subsample the point clouds.
We sample 750 points from each wing silhouette using the farthest-point sampling method (which ensures good coverage of the shape), then compute 1-dimensional Rips persistence:

\textbf{********* Figure : 3.1 -- 70-element Vector{Matrix{Float64}}}


\subsubsection{Strategy 2: Radial filtration}\quad


The \textbf{radial filtration} assigns to each foreground pixel a value equal to its distance from the centroid of the wing. 
Sublevel-set persistence on this function captures how topological features are distributed from the center of the wing outward. 
This is complementary to the Rips filtration, which captures global loop structure without spatial information.

We compute \(H_0\) persistence for the radial filtration, capturing how disconnected vein segments merge as the radial sweep grows outward.



\textbf{***** Figure - 3.4-Radial Filtration of 1 wing}


\subsubsection{Strategy 3: Directional (height) filtration}
\quad

 
The \textbf{directional (height) filtration} sweeps a hyperplane across the wing along a given direction. 
Each foreground pixel is assigned a value equal to its projection onto the sweep direction. 
As the sweep progresses, disconnected vein segments merge — captured by \(H_0\) persistence. 
We use four directions: horizontal \([1,0]\), vertical \([0,1]\), and both diagonals \([1,1]\) and \([1,-1]\).
 

\textbf{**** Figure Directional filtration - of 1 wing}














\section{Topological feature extraction}\quad

For the \textbf{Vietoris-Rips filtration} on connected point clouds, \(H_0\) is uninformative (single infinite bar), so we use only \(H_1\). 
For the \textbf{radial filtration} (computed via sublevel set persistence on the pixel grid), we use only \(H_0\), which captures when disconnected vein segments merge as the radial sweep grows outward — directly encoding vein count and branching patterns. 
Some features are disconsidered, for example the radial \(H_1\) is omitted because pixelate binary images produce very few clean loops under this filtration. 
We also compute \textbf{directional (height) \(H_0\) persistence} in four directions, the axes coordinates and the two diagonals, capturing how vein connectivity changes along different directions.


The persistent homology, given a shape or dataset, tracks how topological features (connected components, loops, voids, etc) appear and disappear as we “grow” the shape through a filtration parameter. 
Each feature has a \textbf{birth} time (when it appears) and a \textbf{death} time (when it gets filled in). 
The collection of all pairs (birth, death) is called a \textbf{persistence diagram}.  Following, some examples of persistence diagram with different filtration are presented: 

\noindent \textbf{COLOCAR Examples: persistence diagrams from each strategy}


\subsection{Summary statistics}\quad

Features with long lifetimes in the persistence diagram represent genuine topological structure, while short lifetimes are typically noises.
We extract extended summary statistics from each persistence diagram using 
\linebreak\textbf{pd\_statistics\_extended}, then remove skewness and kurtosis, that are near to zero relevance in tree based models. We consider the following:
{\bf
\begin{itemize}

\item Count of intervals, max/total/total² persistence

\item Quantiles (10th, 25th, 50th, 75th, 90th)

\item Entropy, std of persistence

\item Median birth, median death, std birth, std death

\item Mean midlife = mean of (birth + death)/2

\item Persistence range = max - min persistence
\end{itemize}}

% So, was dropped stats kurtosis, skewness
After kurtosis and skewness were dropped, 17 features in the statistics per diagram were retained.

% Statistics per diagram retained and we use 17 features from the diagrams.

  Rips \(H_1\): (70, 17)
  
  Radial \(H_0\): (70, 17)
  
  Dir \(H_0\) (horiz): (70, 17)
  
  Dir \(H_0\) (vert): (70, 17)
  
  Dir \(H_0\) (diag1): (70, 17)
  
  Dir \(H_0\) (diag2): (70, 17)

Follow some examples of this features of the diagrams:

\noindent \textbf{COLOCAR Examples: Statistics comparison by family}


\section{Statistics Classification using features}\quad

In this section, we present some statistical methods that, using features, classify images with varying accuracy. Some methods obtain considerable accuracy, while others do not. The decision tree helps to identify which features were most important in obtaining the results.
We build the features matrix from the extended summary statistics of all filtrations:

\textbf{**** Figure - Feature matrix} 

(Feature matrix: (70, 102)

Retained stats per diagram: 17

102 features × 70 samples

Feature-to-sample ratio: 1.46)

We use the leave-one-out cross-validation (LOOCV) with only 72 samples and for each sample, the classifier is trained on all other samples and tested on the held-out one. LOOCV has low bias and is the standard validation strategy for small datasets.

\subsection{Decision tree}\quad

We use a single decision tree with 8 depth as our most interpretable classifier. The tree structure itself provides readable
classification rules:


\textbf{******  Figures - 4.1 Data Frames + Decision trees}


\subsection{Linear Discriminant Analysis (LDA)}
\quad

\textbf{Linear Discriminant Analysis (LDA)} finds a linear projection of the feature space that maximizes the ratio of between class variance to within class variance. 
The projected data is then classified with a simple 1-NN rule.
LDA is a classical method that works well when classes are approximately Gaussian and the number of features is not too large relative to the number of samples.


\textbf{**** Figure - 4.2}

LDA LOOCV: 42/70 (60.0\%)

Balanced accuracy: 55.9\%

Macro-F1: 54.7\%



\subsection{Support Vector Machine (SVM)}\quad

An \textbf{Support Vector Machine (SVM)} finds the hyperplane that maximizes the margin between classes. 
The \textbf{Radial Basis Function (RBF)} \textbf{kernel} maps data into a high-dimensional space where linear separation becomes possible, controlled by a regularization parameter \(C\) (penalty for misclassifications): small \(C\) allows wider margins with more misclassifications, large \(C\) enforces tight boundaries. 
The \textbf{linear kernel} finds a separating hyperplane directly in the original feature space and is less prone to overfitting when \(p>>n\).


\subsubsection{SVM using only Vietoris-Rips \(H_1\) statistics}
\quad

This analysis keeps only features from the 1D Vietoris-Rips barcode (\emph{stats\_rips}) and tests whether loop-based global topology alone is sufficient for classification.

\textbf{**** Figures - 4.6}



% -- SVM with PCA dimensionality reduction (SVM + PCA (Nested LOOCV))
\subsubsection{SVM with PCA dimensionality reduction}\quad

To reduce overfitting risk, we project features with PCA inside each LOOCV fold, then train SVM on the reduced space. 
We scan different variance-retention targets and kernels/costs to identify a good dimensionality.

\textbf{***** Figures 4.5}


% -- SVM (Nested LOOCV)
\subsubsection{SVM (Nested LOOCV)}\quad

Nested cross-validation is a technique used to evaluate the performance of models while tuning their hyperparameters. It helps to reduce bias in model evaluation by using two layers of cross-validation: one for model selection and another for hyperparameter tuning.


\textbf{***** Figures - 4.4}

\subsection{k-NN on Rips Wasserstein distance }\quad

The \textbf{Wasserstein distance} \(W_q\) between two persistence diagrams is the cost of the optimal matching between their points (including matching points to the diagonal, representing trivial features). 
With \(q=1\), it equals the Earth Mover’s Distance; with \( q=2 \), it penalizes large mismatches more heavily.

\textbf{k-Nearest Neighbors (k-NN) }classifies a query by majority vote among its \(k\) nearest neighbors in the distance matrix. 
With \(k=1\), this is the simplest possible classifier, completely hyperparameter-free, and serves as a useful baseline.

As a complementary approach, we compute pairwise Wasserstein distances between the Rips \(H_1\) persistence diagrams and classify with k-NN. 
Unlike the feature-based classifiers above, this operates directly on the persistence diagrams without extracting summary statistics, and is therefore less susceptible to information loss during featurization.

\textbf{***** Figures - 4.7}




% -- -- Balanced Random Forest (Balanced RF (T=200))

\subsection{Balanced Random Forest}\quad

A \textbf{Random Forest} is an ensemble of decision trees, each trained on a bootstrap sample of the data using a random subset of features. 
The final prediction is the majority vote across all trees. \textbf{Balanced} Random Forests oversample minority classes (or weight them inversely to their frequency) so that rare families are not drowned out by common ones. It is important here because Tipulidae has 12 samples while some families have only 2–3. 
Random Forests are robust to overfitting, handle high-dimensional features well, and provide built-in feature importance estimates.

\textbf{****** Figures - 4.3}




\subsection{Top-2 accuracy}\quad 

A single-label classifier must pick exactly one family. Here we relax this to \textbf{top-2}: the model returns the two most probable families and the prediction is considered correct if the true family is among them. This is a useful diagnostic for borderline cases where two families look very similar.

We re-run LOOCV for the best balanced RF, collecting class probabilities via \emph{apply\_forest\_proba}.

\textbf{***** Figures - 5}












%%%%%%%%%%%%%%%%%%%%%%%%%%%%%%%%%%%%%%%%%%%%%%%%%%%%%%%%
%%%%%%%%%%%%%%%%%%%%%%%%%%%%%%%%%%%%%%%%%%%%%%%%%%%%%%%%
\section{Which features drive the classification?}
%%%%%%%%%%%%%%%%%%%%%%%%%%%%%%%%%%%%%%%%%%%%%%%%%%%%%%%%
%%%%%%%%%%%%%%%%%%%%%%%%%%%%%%%%%%%%%%%%%%%%%%%%%%%%%%%%

\subsection{Feature Ablation}

{\bf colocar figura 7}


\subsection{Honest evaluation}\quad


Standard LOOCV can give optimistically biased estimates when hyperparameters are tuned on the same data. \textbf{Nested LOOCV} adds an inner cross-validation loop: for each held-out test sample, the best hyperparameters are selected using only the training fold. This provides an unbiased estimate of generalization performance.

{\bf colocar figura 8}

\subsection{Confusion Matrix}\quad 

{\bf colocar figura 8.1}

\subsection{Misclassified wings: visual inspection}\quad

For each misclassified specimen we show the wing image and radial filtration alongside the 4 nearest neighbors (by Euclidean distance in the z-scored feature space). This helps understand whether errors are due to genuine visual ambiguity or artefacts of segmentation.

{\bf colocar figura 8.2 e imagens exemplos}


















%%%%%%%%%%%%%%%%%%%%%%%%%%%%%%%%%%%%%%%%%%%%%%%%%%%%%%%%
%%%%%%%%%%%%%%%%%%%%%%%%%%%%%%%%%%%%%%%%%%%%%%%%%%%%%%%%
\section{Conclusion and future discussion}\quad
%%%%%%%%%%%%%%%%%%%%%%%%%%%%%%%%%%%%%%%%%%%%%%%%%%%%%%%%
%%%%%%%%%%%%%%%%%%%%%%%%%%%%%%%%%%%%%%%%%%%%%%%%%%%%%%%%

We applied three TDA filtration strategies: Vietoris-Rips, radial, and directional; to classify Diptera families from wing venation images, extracting extended summary statistics per persistence diagram and removing low-relevance moments (skewness and kurtosis).

\subsubsection{Key findings}\quad

\begin{enumerate}
    \item  \textbf{Three filtrations capture complementary information}: The Vietoris-Rips filtration on point-cloud samples captures the global loop structure of the wing venation (number and prominence of wing cells). The radial \(H_0\) filtration captures the center-to-periphery organization: how vein segments merge as the filtration grows outward from the centroid. The directional \(H_0\) filtrations capture vein connectivity along four axes (horizontal, vertical, and both diagonals), encoding how disconnected components merge under directional sweeps.

    \item \textbf{Pruned extended statistics remain sufficient}: After removing skewness and kurtosis, we retain 17 statistics per diagram (count, max/total persistence, quantiles, entropy, median birth/death, etc.). With 6 diagrams × 17 features = 102 total features for 72 samples, the feature-to-sample ratio is ~1.42:1, reducing overfitting risk while preserving discriminative signal.

    \item \textbf{Feature ablation reveals which filtrations matter}: The ablation study compares Rips alone, radial \(H_0\) alone, directional \(H_0\) alone, and their combinations. This provides evidence about whether global topology (Rips), radial organization (radial \(H_0\)), or directional sweep information (directional \(H_0\)) is most discriminative.

    \item Why some filtrations were dropped:
    \begin{list}{$\circ$}{} 
        \item \textbf{Radial \(H_1\)}: On pixelated binary images, few genuine 1-cycles survive the radial filtration, making radial \(H_1\) mostly noise.
        \item \textbf{EDT (Euclidean Distance Transform)}: On binarized images, the EDT is trivially related to the binary structure, providing little additional information beyond what Rips already captures.
        \item \textbf{Cubical (grayscale sublevel-set)}: After binarization, the grayscale information is lost, so cubical persistence reduces to computing persistence on a binary image — equivalent to connected-component analysis.
    \end{list}
    \item \textbf{Nested LOOCV provides honest evaluation}: Standard LOOCV can be optimistic when hyperparameters are tuned on the same data. Nested LOOCV (with 4-fold inner CV for hyperparameter selection) gives unbiased accuracy estimates.

    \item \textbf{Statistical rigor}: We report LOOCV accuracy with Wilson confidence intervals, and nested LOOCV for unbiased evaluation.
\end{enumerate}

\subsubsection{Limitations}\quad

\begin{itemize}
    \item \textbf{Class imbalance}: Tipulidae has 12 samples while some families have only 2–3, which affects classifier performance.
    \item Image quality and preprocessing parameters (blur, threshold) influence topological features.
    With only 72 samples, confidence intervals remain wide regardless of method.
    \item Wings are manually segmented and binarized; automated segmentation could introduce different error patterns.
\end{itemize}

\subsubsection{Future work}\quad

\begin{itemize}
    \item Extend dataset with more specimens per family, especially underrepresented ones
    \item Improve imaging/segmentation quality to reduce noise
    \item Apply extended persistence or zigzag persistence for richer invariants
    \item Investigate which specific topological features (e.g., how many loops, persistence of largest features) correspond to known vein characters in Diptera taxonomy
    \item Try the analysis on non-binarized (grayscale) images, where EDT and cubical filtrations would be more informative
    \item Explore directional \(H_1\) persistence on higher-quality images where loop detection under sweeps may be more reliable
\end{itemize}






\begin{thebibliography}{1}

\bibitem{volovikov1983} A.Yu. Volovikov, \textit{Mappings of free $\mathbb{Z}_{p}$-spaces to manifolds}, Izv. Akad. Nauk SSSR Ser. Mat. \textbf{46} (1) (1982), 36--55; English transl., Math. USSR-Izv. \textbf{20} (1983), 35--53.

\bibitem{volovikov2005} A.Yu. Volovikov, \textit{Coincidence points of mappings $\mathbb{Z}_{p}^{n}$-spaces}, Izv. Ross. Akad.
Nauk Ser. Mat. \textbf{69} (2005), 53--106; transl. in Izv. Math. \textbf{69} (2005), 913--962.

\end{thebibliography}

% ------------------------------------------------------------------------
\end{document}
% ------------------------------------------------------------------------
